\section{Исследуйте на непрерывность: найдите точки разрыва и укажите
точки устранимого разрыва функции двух переменных}
\[
z = \frac{\sin^2{x}\sin{y}}{\sin^4{x} + \sin^2{y}}
\]

\subsection{Область определения}

Дробь определена тогда и только тогда, когда знаменатель не равен нулю:
\[
\sin^4{x} + \sin^2{y} \neq 0.
\]
Так как \(\sin^4{x} \ge 0\) и \(\sin^2{y} \ge 0\),сумма равна нулю тогда и только тогда, когда одновременно \(\sin{x} = 0\) и \(\sin{y} = 0\).Поэтому область определения:
\[
D = \{ (x,y) \in \mathbb{R}^2:(\sin{x},\sin{y}) \neq (0,0) \},
\]
эквивалентно:
\[
D = \mathbb{R}^2 \setminus \{(k\pi, m\pi) \mid k,m \in \mathbb{Z}\}\}.
\]

\subsection{Непрерывность на области определения}

На \( D \) числитель и знаменатель являются непрерывными функциями, кроме того на \( D \) знаменатель строго положителен. Следовательно дробь --- отношение непрерывных функций с ненулевым знаменателем, значит \( z \) непрерывна во всех точках \( D \).

\subsection{Анализ накопленных точек \((k\pi, m\pi)\)}

Исследуем предел функции при \((x, y) \xrightarrow{} (k\pi, m\pi)\).Пусть \(\sin{y} = c\sin^2{x}\):
\[
z = \frac{c}{1 + c^2}.
\]
Значение предела зависит от параметра \(c\).
Так как пределы по разным путям различаются, общий предел при \((x, y) \xrightarrow{} (k\pi, m\pi)\) не существует.

\subsection{Итоговый результат}

\begin{itemize}
    \item Функция непрерывна в \(D = \mathbb{R}^2 \setminus \{(k\pi, m\pi) \mid k,m \in \mathbb{Z}\}\}\).
    \item Точки разрыва:\((k\pi, m\pi) ,\ k,m \in \mathbb{Z}\).
    \item Точки устранимого разрыва не существуют.
\end{itemize}
